\documentclass{jsarticle}
\usepackage[dvipdfmx]{graphicx}
\usepackage[dvipdfmx]{hyperref}
\usepackage[section]{placeins}
\usepackage{here}
\usepackage{amsmath}
\usepackage{amssymb}
\usepackage{amsfonts}
\usepackage{listings,jvlisting}
\usepackage{url}
\usepackage{float}
\lstset{
  basicstyle={\ttfamily},
  identifierstyle={\small},
  commentstyle={\smallitshape},
  keywordstyle={\small\bfseries},
  ndkeywordstyle={\small},
  stringstyle={\small\ttfamily},
  frame={tb},
  breaklines=true,
  columns=[l]{fullflexible},
  numbers=left,
  xrightmargin=0zw,
  xleftmargin=3zw,
  numberstyle={\scriptsize},
  stepnumber=1,
  numbersep=1zw,
  lineskip=-0.5ex
}

\title{進捗報告}
\author{冨田　哲志}
\date{2026/3/10}

\begin{document}

\maketitle

\section{先週の進捗}
\begin{itemize}
    \item 研修課題大問3
     \begin{itemize}
         \item データの取り直し
         \item Dropoutの導入
         \item Action Chunkingの導入
     \end{itemize}
    
        
\end{itemize}

\section{研修課題大問3}
\subsection{データの取り直し}
これまでスタートの位置を真ん中にした30個のデータでうまく推論できていなかったため、
\subsection{関節データにDropoutの導入}
\subsection{Action Chunkingの導入}




\subsection{}

% \begin{figure}[H] % [htbp] は画像を出す位置の指定
%   \centering         % 中央揃え
%   \includegraphics[width=10cm]{} % 画像ファイル名とサイズ
%   \caption{元画像と再構成画像の比較} % 図の下に表示される説明
%   \label{fig:sample}   % 本文中で参照するためのラベル
% \end{figure}









\section{今週の予定}
\begin{itemize}
    \item 研修課題大問3
\end{itemize}

\end{document}